% !TeX root = ../初等数论笔记.tex
% 由主文档通过 \input 载入；请编译主文档。
\setcounter{chapter}{1} % 保留原书章号。
\chapter{习题选答与综合训练}
\label{app:exercises}

\section{各章计算题选答}

\textbf{第\ref{ch:introduction}章。}相邻数问题可将 $1,\ldots,2n$ 分为 $n$ 对 $\{1,2\},\{3,4\},\ldots,\{2n-1,2n\}$ 后应用抽屉原理。证明 $\sqrt3$ 无理时，先按模 $3$ 的余数 $0,1,2$ 分类，得到平方被 $3$ 整除则原数被 $3$ 整除，再重复最小分母论证。

\textbf{第\ref{ch:divisibility}章。}$-37=8(-5)+3$；$\gcd(414,662)=2=8\cdot414-5\cdot662$。$35x+22y=1$ 的全部解为 $x=-5+22t,y=8-35t$（$t\in\mathbb Z$）。

\textbf{第\ref{ch:primes}章。}$v_2(100!)=50+25+12+6+3+1=97$；$v_5\bigl(\binom{100}{50}\bigr)=24-2\cdot12=0$。对 $1\le k<p$，$\binom pk$ 的分子含一个因子 $p$，分母 $k!(p-k)!$ 不含 $p$，故 $p\mid\binom pk$。

\textbf{第\ref{ch:congruences}章。}$17^{-1}\equiv38\pmod{43}$，因为 $17\cdot38=15\cdot43+1$。平方模 $8$ 的全部可能余数为 $0,1,4$。

\textbf{第\ref{ch:crt}章。}$14x\equiv8\pmod{30}$ 的解为 $x\equiv7,22\pmod{30}$。两条相容约束的解为 $x\equiv13\pmod{24}$，把余数 $5$ 改为 $4$ 后奇偶性冲突而无解。模 $35$ 的平方根 $1$ 有四个：$1,6,29,34$。

\textbf{第\ref{ch:arithmetic-functions}章。}$840=2^3\cdot3\cdot5\cdot7$，故 $\varphi(840)=192,\tau(840)=32,\mu(840)=0$。卷积恒等式由 $\tau=\mathbf1*\mathbf1$ 与 $\mu*\mathbf1=\varepsilon$ 得到 $\mu*\tau=\mathbf1$。

\textbf{第\ref{ch:classical-congruences}章。}$7^{222}\equiv12\pmod{13}$、$3^{100}\equiv1\pmod{20}$、$5^{117}\equiv1\pmod{19}$。若 $p>3$ 为素数，$p$ 为奇数给出 $p^2\equiv1\pmod8$，$3\nmid p$ 给出 $p^2\equiv1\pmod3$，再合并为模 $24$ 的结论。

\textbf{第\ref{ch:primitive-roots}章。}模 $7$ 的原根为 $3,5$；$10$ 模 $13$ 的阶为 $6$。$x^3\equiv1\pmod7$ 的根为 $1,2,4$；$x^2\equiv2\pmod{49}$ 的根为 $10,39$。若 $g$ 是模 $m$ 的原根，由式\eqref{eq:order-power}，$g^k$ 为原根当且仅当 $\gcd(k,\varphi(m))=1$，所以原根类有 $\varphi(\varphi(m))$ 个。

\textbf{第\ref{ch:quadratic-residues}章。}模 $13$ 的非零平方剩余为 $1,3,4,9,10,12$；$\left(\frac7{19}\right)=1$、$\left(\frac{13}{17}\right)=1$；$x^2\equiv4\pmod{45}$ 的根为 $2,7,38,43$。

\textbf{第\ref{ch:diophantine}章。}斜边为 $25$ 的三元组为 $(7,24,25)$、$(15,20,25)$。$45=3^2+6^2$、$130=3^2+11^2$，而 $75$ 不可表示，因为其素因子 $3$ 的指数为奇数。方程 $xy=2x+3y$ 化为 $(x-3)(y-2)=6$；由 $x(y-2)=3y>0$ 知 $y>2$，继而 $x>3$，所以只有正因子对，得到 $(4,8),(5,5),(6,4),(9,3)$。

\textbf{第\ref{ch:continued-fractions}章。}$415/93$ 的渐近分数为 $4,9/2,58/13,415/93$。$\sqrt2$ 的前五个渐近分数为 $1,3/2,7/5,17/12,41/29$；$\sqrt3=[1;\overline{1,2}]$、$\sqrt{13}=[3;\overline{1,1,1,1,6}]$。

\textbf{第\ref{ch:pell}章。}$D=3,5,13$ 的基本解分别为 $(2,1),(9,4),(649,180)$。负方程 $x^2-2y^2=-1$ 的前三组正解为 $(1,1),(7,5),(41,29)$。平方三角数的前两个正指标为 $n=1,8$。

\textbf{第\ref{ch:prime-distribution}章。}$6!+2,\ldots,6!+6$ 给出连续五个合数。仅由集合无穷不能推断倒数和发散，例如 $\{2^k:k\ge1\}$ 的倒数和收敛；要讨论某个素数子集，仍需关于该子集分布的额外信息。

\textbf{第\ref{ch:algorithms}章。}$13^{27}\equiv7\pmod{55}$。对 $561$ 取底数 $2$，$560=2^4\cdot35$，相继得到 $263,166,67,1$；起点非 $1$，且此前未出现 $-1\equiv560$，所以未通过 Miller--Rabin 检验。$10403=102^2-1^2=101\cdot103$。

\section{分层习题提示与方法讨论}

本节按章节给出后四题的关键步骤。计算题附结果，证明题说明主要转化；使用时应尝试补全中间推导，并回到正文检查每个所用结论的条件。习题编号由各章独立计数，以下同时注明章节。

\textbf{第\ref{ch:introduction}章。}习题\ref{ex:intro-negation}的否定为“对每个整数 $n$，存在整数 $m$ 使 $m>n$”，取 $m=n+1$ 即可；量词由“存在、任意”依次变为“任意、存在”，不等号也要同时否定。习题\ref{ex:intro-strong}先验证 $2=2\cdot1+3\cdot0$、$3=2\cdot0+3\cdot1$；对 $n\ge4$，由 $n-2\ge2$ 的表示加上一个 $2$ 得到 $n$ 的表示。两个起点分别覆盖之后的偶数与奇数。

习题\ref{ex:intro-dividing-pair}把每个所选数唯一写成 $2^ru$，其中 $u$ 为奇数。区间中可能的奇数部分只有 $1,3,\ldots,2n-1$ 这 $n$ 个；抽屉原理保证两个数具有相同的 $u$，指数较小的数就整除较大的数。习题\ref{ex:intro-descent-domain}中，正实数集合不具有正整数的良序性质；给出的值集没有最小元素。因此递降时不仅要证明参数变小，还要证明它始终属于正整数或另一个不能无限严格下降的集合。

\textbf{第\ref{ch:divisibility}章。}习题\ref{ex:div-consecutive}使用 $(n+1)-n=1$，习题\ref{ex:div-combinations}使用 $3(2n+1)-2(3n+1)=1$。任意公因数都整除右侧的 $1$，故最大公因数为 $1$；这比先分别分解两个随 $n$ 变化的整数更直接。习题\ref{ex:div-nonnegative}中，模 $4$ 得 $y$ 被 $4$ 整除，写 $y=4t,x=25-7t$，非负性给出 $t=0,1,2,3$，所以解为 $(25,0),(18,4),(11,8),(4,12)$。

习题\ref{ex:div-gcd-lcm-count}在 $d\nmid L$ 时无解。若 $d\mid L$，写 $a=du,b=dv$，则 $\gcd(u,v)=1$、$uv=L/d$。由唯一分解，$L/d$ 的每个不同素因子的整个幂次必须全部分给 $u$ 或全部分给 $v$，不能拆给两边。若不同素因子数为 $r$，有序对共有 $2^r$ 个；$L=d$ 时 $r=0$，唯一解为 $(d,d)$。

\textbf{第\ref{ch:primes}章。}习题\ref{ex:prime-cube}中 $180=2^2\cdot3^2\cdot5$，补足到 $3$ 的倍数所需的最小因子为 $2\cdot3\cdot5^2=150$，此时 $180\cdot150=30^3$。习题\ref{ex:prime-odd-divisors}把约数 $d$ 与 $n/d$ 配对；只有 $d^2=n$ 才会出现一个不与不同约数配对的对象。指数证明则使用约数个数 $\prod_i(\alpha_i+1)$，它为奇数当且仅当每个 $\alpha_i$ 都为偶数。

习题\ref{ex:prime-rational-root}设 $\sqrt n=a/b$，其中 $a,b>0$ 互素。若 $b>1$，取其素因子 $p$，由 $a^2=nb^2$ 得 $p\mid a$，与互素矛盾；所以 $b=1$。习题\ref{ex:prime-binomial}使用阶乘指数：$v_p((p^2)!)=p+1$、$v_p(p!)=1$、$v_p((p^2-p)!)=p-1$，相减得 $1$。在 $p=2$ 时计算仍成立，例如 $\binom42=6$ 恰含一个因子 $2$。

\textbf{第\ref{ch:congruences}章。}习题\ref{ex:cong-cancel}约分后应为 $x\equiv1\pmod3$，在模 $12$ 下为 $1,4,7,10$ 四类；错误保留模数只会留下 $1$。习题\ref{ex:cong-obstruction}中，平方模 $3$ 为 $0$ 或 $1$，所以平方加一不被 $3$ 整除；但 $x=2,y=1$ 已构成第二个方程的整数解，说明无解障碍依赖具体模数。

习题\ref{ex:cong-last-digit}用模 $10$ 的周期 $7,9,3,1$，因 $2026$ 除以 $4$ 余 $2$，个位为 $9$。习题\ref{ex:cong-palindrome}设数位为 $a_0,\ldots,a_{2r-1}$，回文条件为 $a_j=a_{2r-1-j}$；在交错数位和中，这两个位置的指数奇偶相反，贡献相互抵消，因此总和为零，原整数被 $11$ 整除。

\textbf{第\ref{ch:crt}章。}习题\ref{ex:crt-redundant}中 $x\equiv6\pmod8$ 已蕴含 $x\equiv2\pmod4$，所以最终解仍是 $x\equiv6\pmod8$。习题\ref{ex:crt-three}先写 $x=2+6t$，由第二式得 $6t\equiv3\pmod9$，即 $t\equiv2\pmod3$，故 $x\equiv14\pmod{18}$。再写 $x=14+18s$，模 $5$ 得 $s\equiv0\pmod5$，最终为 $x\equiv14\pmod{90}$。

习题\ref{ex:crt-zero-roots}要求 $2\mid x$ 且 $3\mid x$，即 $6\mid x$，所以模 $12$ 有 $0,6$ 两根。习题\ref{ex:crt-idempotents}在每个素数模数下有 $x(x-1)\equiv0$，因此恰有 $0,1$ 两个局部选择；中国剩余定理给出 $2^r$ 个组合。模 $30$ 的全部根为 $0,1,6,10,15,16,21,25$。

\textbf{第\ref{ch:arithmetic-functions}章。}习题\ref{ex:arith-noncoprime}的值为 $4,6,72$；因为 $12$ 与 $18$ 不互素，积性不适用，且 $72\ne4\cdot6$。习题\ref{ex:arith-invert-twelve}得到 $f(12)=F(12)-F(6)-F(4)+F(2)$，来自 $d=1,2,3,6$ 的莫比乌斯值；含平方因子的 $d=4,12$ 不贡献。习题\ref{ex:arith-squarefree}的结果为 $20-5-2=13$，对应 $1,2,3,5,6,7,10,11,13,14,15,17,19$；$d=4$ 的莫比乌斯值为零。

习题\ref{ex:arith-inverse}先在 $n=1$ 处得到 $f(1)g(1)=1$，故 $f(1)\ne0$ 为必要条件。若此条件满足，令 $g(1)=1/f(1)$，再按 $n>1$ 递推
\[
 g(n)=-\frac1{f(1)}\sum_{\substack{d\mid n\\d>1}}f(d)g(n/d).
\]
求和中 $n/d<n$，所以右侧已经确定。该定义使每个 $n>1$ 的卷积值为零，证明存在性；任何逆元都必须满足相同递推，又证明唯一性。

\textbf{第\ref{ch:classical-congruences}章。}习题\ref{ex:classical-noncoprime}中，$6^2\equiv1\pmod{35}$，故第一个余数为 $1$。模 $40$ 时拆成模 $8$ 与模 $5$：$6^{100}$ 分别余 $0$ 与 $1$，合并得到 $16$。习题\ref{ex:classical-factorial}中 $12!\equiv-1\pmod{13}$，而 $11\cdot12\equiv2$，故 $2\cdot10!\equiv-1$，乘以 $2^{-1}\equiv7$ 得 $10!\equiv6$。

习题\ref{ex:classical-fermat-all}分别证明 $2,3,5$ 整除 $a^5-a$。模 $5$ 直接用费马小定理的全整数形式；模 $3$ 若 $3\mid a$ 则两项为零，否则 $a^2\equiv1$，故 $a^5\equiv a$；模 $2$ 同理。三个模数两两互素，所以乘积 $30$ 整除差。习题\ref{ex:classical-pseudoprime}中底数 $2$ 的结论见正文；底数 $3$ 则有 $3^{340}\equiv3^{10}\equiv25\not\equiv1\pmod{31}$，因此更不可能模 $341$ 余 $1$。这里 $340\equiv10\pmod{30}$，使用的是模 $31$ 的费马周期。

\textbf{第\ref{ch:primitive-roots}章。}习题\ref{ex:order-list}对类 $1,2,4,5,7,8$ 求得的阶依次为 $1,6,3,6,3,2$，所以原根为 $2,5$。习题\ref{ex:order-composite}中，$a^k\equiv1\pmod{mn}$ 当且仅当它模 $m$、模 $n$ 都为 $1$，也就是两个局部阶都整除 $k$，最小这样的正整数即为最小公倍数。

习题\ref{ex:order-sixth}取模 $13$ 的原根 $2$，写 $x=2^t$，则 $6t\equiv0\pmod{12}$ 等价于 $t$ 为偶数，得到根 $1,3,4,9,10,12$。习题\ref{ex:order-singular}中，模 $9$ 的第一式有 $0,3,6$ 三根，第二式无根；而 $x^2\equiv0\pmod{27}$ 要求 $2v_3(x)\ge3$，等价于 $9\mid x$，所以根为 $0,9,18$。$x=0$ 单独满足方程，赋值记号本身只用于非零整数。

\textbf{第\ref{ch:quadratic-residues}章。}习题\ref{ex:quadratic-sign}有 $\left(\frac3{17}\right)=\left(\frac{17}3\right)=\left(\frac23\right)=-1$、$\left(\frac5{29}\right)=\left(\frac{29}5\right)=\left(\frac45\right)=1$；两次交换都不变号，因为至少一个素数模 $4$ 余 $1$。习题\ref{ex:quadratic-jacobi}有 $\left(\frac53\right)=\left(\frac57\right)=-1$，所以 $\left(\frac5{21}\right)=1$，但两个局部方程都无解，整体也无解。

习题\ref{ex:quadratic-two-power}可逐个检验奇数代表元，得到 $3,13,19,29$，恰有四根；也可先取 $3$，再利用取负与加 $16$ 得到另外三根。习题\ref{ex:quadratic-valuation}先由模 $3$ 条件得 $x=3y$，约去 $9$ 后要求 $y^2\equiv1\pmod3$，即 $3\nmid y$。由于 $x$ 在模 $27$ 下计数，$y$ 应在模 $9$ 下计数，共有 $1,2,4,5,7,8$ 六类；因此 $x=3,6,12,15,21,24$。

\textbf{第\ref{ch:diophantine}章。}习题\ref{ex:dio-factors}写成 $(x-y)(x+y)=45$。因 $x,y\ge0$，有 $x>y$，两个正因子应同奇且前者不大于后者；因子对为 $(1,45),(3,15),(5,9)$，得到 $(23,22),(9,6),(7,2)$。习题\ref{ex:dio-hypotenuse}中，本原斜边为奇数。若其某个素因子 $p\equiv3\pmod4$，则 $p\mid x^2+y^2$，由推论\ref{cor:prime-three-mod-four}迫使 $p\mid x,y$，与本原性矛盾。

习题\ref{ex:dio-two-square-list}中 $50=1^2+7^2$、$98=7^2+7^2$、$325=1^2+18^2$；$77=7\cdot11$ 的两个模 $4$ 余 $3$ 的素因子均为奇数指数，故不可表示。习题\ref{ex:dio-area}中，本原三角形的面积为 $uv(u^2-v^2)$。因为 $u,v$ 奇偶不同，面积被 $2$ 整除；若 $3\mid uv$ 则也被 $3$ 整除，否则 $u^2\equiv v^2\equiv1\pmod3$，差被 $3$ 整除。一般三角形的面积为本原面积乘以一个整数平方，因此仍被 $6$ 整除。

\textbf{第\ref{ch:continued-fractions}章。}习题\ref{ex:cf-negative}的答案为 $[-3;1,2]$，因为下取整是 $-3$，而不是向零截断为 $-2$。习题\ref{ex:cf-sqrt-seven}中，状态 $(B_j,C_j)$ 从 $(0,1)$ 起依次为 $(2,3),(1,2),(1,3),(2,1)$，之后回到 $(2,3)$，故 $\sqrt7=[2;\overline{1,1,1,4}]$。前四个渐近分数为 $2,3,5/2,8/3$，平方后与 $7$ 比较即可核对交错方向。

习题\ref{ex:cf-neighbors}由严格夹逼知 $cv-du$ 与 $bu-av$ 都是正整数。利用 $bc-ad=1$，有
\[
 v=b(cv-du)+d(bu-av)\ge b+d.
\]
这说明相邻行列式为 $1$ 的分数之间，要插入一个分数就必须使用较大的分母。习题\ref{ex:cf-golden}的周期尾部与原式相同，故 $\alpha=1+1/\alpha$，满足 $\alpha^2-\alpha-1=0$；由 $\alpha>1$ 得 $\alpha=(1+\sqrt5)/2$。

\textbf{第\ref{ch:pell}章。}习题\ref{ex:pell-seven}的周期长为 $4$，基本解为 $(8,3)$，第二组由 $(8+3\sqrt7)^2=127+48\sqrt7$ 得 $(127,48)$。习题\ref{ex:pell-negative-five}中最小负范数单位为 $2+\sqrt5$，正方程基本单位为 $9+4\sqrt5$，故前三组为 $(2,1),(38,17),(682,305)$；代入分别得到 $-1$。

习题\ref{ex:pell-second-order}令 $U=x_1+y_1\sqrt D$，因 $U^{-1}=x_1-y_1\sqrt D$，有 $U+U^{-1}=2x_1$，从而 $U^{n+2}-2x_1U^{n+1}+U^n=0$。比较有理部分与 $\sqrt D$ 的系数，得到两个递推；比较系数合法是因为 $\sqrt D$ 无理。习题\ref{ex:pell-obstruction}模相应素数 $p$ 得 $x^2\equiv-1\pmod p$，与 $p\equiv3\pmod4$ 矛盾。没有这样的素因子只表示这一局部障碍消失，是否有解仍须使用周期奇偶等充分判据。

\textbf{第\ref{ch:prime-distribution}章。}习题\ref{ex:distribution-sparse}可取 $\{2^j:j\ge1\}$，其倒数和为几何级数 $\sum_{j\ge1}2^{-j}=1$。习题\ref{ex:distribution-six-five}若假设素数只有有限个 $p_1,\ldots,p_r$，令 $N=6p_1\cdots p_r-1$。它与 $6$ 互素，所有素因子都模 $6$ 余 $1$ 或 $5$；若没有余 $5$ 的因子，乘积只能余 $1$。因此存在一个余 $5$ 的新素因子，与有限列表矛盾。

习题\ref{ex:distribution-density}直接得到 $\pi(x)/x\sim1/\log x\to0$。自然密度为零只控制累计个数相对于 $x$ 的比例，并不保证倒数和收敛；素数的这两个已知结论正说明两种稀疏性标准不同。习题\ref{ex:distribution-asymptotic}可取 $f(x)=x+\sqrt x$、$g(x)=x$（$x>0$），比值为 $1+1/\sqrt x\to1$，差却为 $\sqrt x\to\infty$。

\textbf{第\ref{ch:algorithms}章。}习题\ref{ex:algo-inverse-fail}若逆元存在，则 $6u-15v=1$ 有整数解，但左侧总被 $3$ 整除，因此不可能；算法应返回“逆元不存在，最大公因数为 $3$”。习题\ref{ex:algo-rsa-small}有 $n=33,\varphi(n)=20$，$3\cdot7\equiv1\pmod{20}$，故 $d=7$。消息 $6$ 加密为 $6^3\bmod33=18$，解密为 $18^7\bmod33=6$，覆盖了消息与模数不互素的情形。

习题\ref{ex:algo-sqrt-factor}若 $\gcd(r-1,n)=1$，由 $n\mid(r-1)(r+1)$ 可消去 $r-1$，推出 $r\equiv-1\pmod n$，矛盾；它又不能等于 $n$，否则 $r\equiv1$。交换两个因子得到另一个最大公因数也在 $1$ 与 $n$ 之间。习题\ref{ex:algo-test-cases}可依次取 $(a,e,m)=(-2,5,7),(5,0,13),(6,4,8),(2,3,0)$，预期输出分别为 $3,1,0$，以及因不满足 $m\ge2$ 而拒绝。有限测试能发现具体错误，而一般正确性仍来自循环不变式与终止性证明。

\section{综合例题}

\begin{example}[同余化简与解数]
求 $6x\equiv9\pmod{15}$、$x\equiv2\pmod7$ 的全部整数解，并说明模 $105$ 有几个解类。
\end{example}
\begin{solve}
第一式等价于 $x\equiv4\pmod5$。与第二式合并，得 $x\equiv9\pmod{35}$。因此模 $105$ 的解类为 $9,44,79$，共三个。模 $35$ 唯一与模 $105$ 有三类是同一整数解集的两种表达。
\end{solve}

\begin{example}[最大公因数的求和]
对 $n\ge1$，证明
\[
 \sum_{k=1}^n\gcd(k,n)=\sum_{d\mid n}d\varphi(n/d).
\]
\end{example}
\begin{proof}
按 $\gcd(k,n)=d$ 分类。写 $k=du$ 后，条件等价于 $1\le u\le n/d$、$\gcd(u,n/d)=1$，故该类含 $\varphi(n/d)$ 个数，每个贡献 $d$。对全部正约数求和即可。例如 $n=12$ 时，右侧为 $4+4+6+8+6+12=40$。
\end{proof}

\begin{example}[平方和的全部表示]
求 $x^2+y^2=2025$ 的全部非负整数解，不区分坐标次序。
\end{example}
\begin{solve}
$2025=3^4\cdot5^2$。因 $3\equiv3\pmod4$，推论\ref{cor:prime-three-mod-four}使 $x,y$ 同时被 $3$ 整除；除以 $9$ 后再次应用，得 $9\mid x,y$。写 $x=9a,y=9b$，则 $a^2+b^2=25$。规定 $a\ge b\ge0$，有 $2b^2\le25$，故只需检查 $b=0,1,2,3$，得到 $(a,b)=(5,0),(4,3)$。原方程只有 $(45,0),(36,27)$。
\end{solve}

\begin{example}[平方三角数]
求使 $n(n+1)/2$ 为平方的全部正整数 $n$ 的生成方法。
\end{example}
\begin{solve}
设 $n(n+1)/2=m^2$，配方得 $(2n+1)^2-8m^2=1$。$D=8$ 的基本解为 $(3,1)$，所以令
\[
 X_k+Y_k\sqrt8=(3+\sqrt8)^k,\qquad
 n_k=\frac{X_k-1}{2},\quad k\ge1.
\]
方程保证 $X_k$ 为奇数，故每个 $n_k$ 为正整数；反向变换和 Pell 方程的完备性保证没有遗漏。前三个指标为 $1,8,49$，对应平方三角数 $1,36,1225$。
\end{solve}

